\documentclass[11pt,a4paper]{article}

% --- Codificación estándar para pdflatex ---
\usepackage[utf8]{inputenc}
\usepackage[T1]{fontenc}

% --- Idioma sin opciones incompatibles con pdflatex ---
\usepackage[spanish,es-noquoting]{babel}

% --- Matemática y símbolos ---
\usepackage{amsmath,amssymb,amsthm}

% --- Geometría de página ---
\usepackage{geometry}
\geometry{margin=2.5cm}

% --- Gráficos y TikZ (con compatibilidad de babel) ---
\usepackage{tikz}
\usetikzlibrary{automata, positioning, arrows.meta, babel}

% --- Tablas y colores ---
\usepackage{booktabs}
\usepackage{xcolor}
\usepackage{colortbl}
\usepackage{enumitem}

% --- Hipervínculos (cargar siempre al final del preámbulo) ---
\usepackage{hyperref}
\hypersetup{
    colorlinks=true,
    linkcolor=blue!80!black,
    urlcolor=blue!80!black
}
\usepackage{bookmark}

\setlist[itemize]{label=--, leftmargin=1.5em, itemsep=2pt, parsep=0pt}
\setlist[enumerate]{leftmargin=1.5em, itemsep=2pt, parsep=0pt}

% Definición de colores sobrios para bloques
\definecolor{primaryblue}{RGB}{24, 43, 73}
\definecolor{accentblue}{RGB}{41, 98, 153}
\definecolor{graybg}{RGB}{246, 248, 250}

\hypersetup{
    colorlinks=true,
    linkcolor=accentblue,
    urlcolor=accentblue
}

\theoremstyle{definition}
\newtheorem{definicion}{Definición}
\newtheorem{proposicion}{Proposición}
\newtheorem{teorema}{Teorema}
\newtheorem{ejemplo}{Ejemplo}
\newtheorem{observacion}{Observación}
\newtheorem{demostracion}{Demostración}
\newtheorem{propiedades}{Propiedades}

\setlength{\parindent}{0pt}
\setlength{\parskip}{5pt}

\begin{document}

\begin{center}
    {\LARGE \textbf{\color{primaryblue} Guía de Estudio: Procesos y Matrices de Markov}}\\[4pt]
    {\large \textbf{Álgebra Lineal Computacional (ALC)}} \\[2pt]
    \small Resumen teórico, herramientas prácticas para ejercicios y conexiones estructurales
\end{center}
\vspace{-0.5em}
\hrule height 1.2pt
\vspace{1em}

\section*{Motivación: El Modelo Google PageRank y Caminatas Aleatorias}

\subsection*{1. Planteo intuitivo: ¿Cómo ordenar la web por relevancia?}
En los inicios de los motores de búsqueda, determinar la importancia relativa
de una página web se convirtió en un desafío central. Una primera
aproximación intuitiva consiste en contar la cantidad de hipervínculos que
apuntan a cada sitio. Sin embargo, no todos los enlaces tienen el
mismo peso: recibir una recomendación o enlace desde una página muy influyente
(como Wikipedia) otorga más relevancia que recibirla desde un sitio
desconocido.

Para resolver esta aparente circularidad (''una página es importante si es
referenciada por páginas importantes''), el algoritmo de Google PageRank modela
la navegación web mediante un \textbf{caminante aleatorio} (\textit{random
    surfer}):
\begin{enumerate}
    \item Imaginemos una población fija de internautas navegando por una red.
    \item En el instante inicial $k = 0$, cada visitante elige una página departida.
    \item En cada paso de tiempo discreto $k \to k+1$, si el navegante se encuentra en una página con $d$ enlaces salientes disponibles, elije al azar cualquiera de ellos de manera equiprobable con probabilidad $1/d$.
    \item Nos preguntamos: tras navegar mucho tiempo ($k \to \infty$), ¿cuál es la probabilidad o proporción de visitantes que se encontrará en cada página
\end{enumerate}

\subsection*{2. Diagrama de Estados y Transiciones}
Consideremos una subred compuesta por 4 sitios web interconectados: $A$,
$B$, $C$ y $D$.
\begin{itemize}
    \item Desde la página $A$ hay enlaces hacia $B$ o puede quedarse en $A$ (2 enlaces disponibles).
    \item Desde la página $B$ hay 3 enlaces salientes, hacia $A$, $C$ o puede quedarse en $B$ (3 enlaces disponibles).
    \item Desde la página $C$ el internauta puede quedarse en $C$ (autoenlace) o pasar a $B$ (2 enlaces disponibles).
    \item Desde la página $D$ hay enlaces salientes hacia $A$ y $C$ y puede quedarse en $D$ (autoenlace) (3 enlaces disponibles).
\end{itemize}

El siguiente diagrama de transiciones ilustra el sistema de estados discretos:

\begin{center}
    \begin{tikzpicture}[
            node distance=3.2cm,
            auto,
            every state/.style={circle, draw=blue!70, fill=blue!10, thick, minimum size=1.1cm, font=\bfseries}
        ]
        % Estados / Páginas
        \node[state] (A) {$A$};
        \node[state] (B) [right of=A] {$B$};
        \node[state] (C) [below of=B] {$C$};
        \node[state] (D) [below of=A] {$D$};

        % Flechas de transición con sus probabilidades
        \path[->, >=stealth, thick]
        (A) edge [bend left=15] node[above] {$1/4$} (B) 
        (A) edge [loop left=15] node {$1/2$} (A)
        (B) edge [bend left=15] node[below] {$1/4$} (A) 
        (B) edge [bend left=15] node[right] {$1/4$} (C)
        (B) edge [loop right=15] node {$1/2$} (B) 
        (C) edge [bend left=15] node[left]  {$1/4$} (B)
        (C) edge [loop right=15] node {$1/2$} (C) 
        (D) edge [bend left=15] node[left]  {$1/4$} (A)
        (D) edge [bend left=15] node[below]  {$1/4$} (C)
        (D) edge [loop left=15] node {$1/2$} (D);
    \end{tikzpicture}
\end{center}

\subsection*{3. Del Grafo a la Matriz de Transición de Markov}
Este proceso se puede plantear matricialmente.

Llamamos $v^{(0)}$ al estado inicial del sistema.

\[
    v^{(0)} = 
        \begin{pmatrix} 
            A{(0)} \\ 
            B^{(0)} \\ 
            C^{(0)} \\ 
            D^{(0)} 
        \end{pmatrix} = 
        \begin{pmatrix} 
            1000 \\ 
            1000 \\ 
            1000 \\ 
            1000 
        \end{pmatrix}
\]

Llamamos $v^{(1)}$ al estado del sistema al minuto. 

\[ 
    v^{(1)} =
    \begin{pmatrix} 
            A{(1)} \\ 
            B^{(1)} \\ 
            C^{(1)} \\ 
            D^{(1)} 
        \end{pmatrix} = 
        \begin{pmatrix}
            1/2 & 1/4 & 0   & 1/4 \\
            1/2 & 1/2 & 1/2 & 0   \\
            0   & 1/4 & 1/2 & 1/4 \\
            0   & 0   & 0   & 1/2
        \end{pmatrix} 
        \begin{pmatrix} 
            A{(0)} \\ 
            B^{(0)} \\ 
            C^{(0)} \\ 
            D^{(0)} 
        \end{pmatrix}
\]

En general

\[  
    v^{(k+1)} =
    \begin{pmatrix} 
            A^{(k+1)} \\ 
            B^{(k+1)} \\ 
            C^{(k+1)} \\ 
            D^{(k+1)} 
        \end{pmatrix} = 
        \begin{pmatrix}
            1/2 & 1/4 & 0   & 1/4 \\
            1/2 & 1/2 & 1/2 & 0   \\
            0   & 1/4 & 1/2 & 1/4 \\
            0   & 0   & 0   & 1/2
        \end{pmatrix} 
        \begin{pmatrix} 
            A^{(k)} \\ 
            B^{(k)} \\ 
            C^{(k)} \\ 
            D^{(k)} 
        \end{pmatrix}
\]

% Para transformar este proceso en un modelo matricial, representamos la
% distribución de navegantes en el instante $k$ mediante un vector de
% probabilidades $v^{(k)} \in \mathbb{R}^4$, donde $v_j^{(k)} \ge 0$ y
% $\sum_{j=1}^4 v_j^{(k)} = 1$:
% \[
%     v^{(k)} = \begin{pmatrix} v_1^{(k)} \\ v_2^{(k)} \\ v_3^{(k)} \\ v_4^{(k)} \end{pmatrix}
% \]

% La regla fundamental de balance de probabilidad indica que la masa de
% probabilidad que llega al estado $i$ proviene de los estados $j$ que lo
% enlazan:
% \[
%     v_i^{(k+1)} = \sum_{j=1}^n P_{ij} \, v_j^{(k)} \implies v^{(k+1)} = P \, v^{(k)}
% \]
% \textbf{Regla de construcción por columnas (Convención estándar en ALC):}
% \begin{itemize}
%     \item La casilla $P_{ij}$ representa la \textbf{probabilidad de saltar del estado $j$ (columna) al estado $i$ (fila)}.
%     \item Por conservación de la probabilidad total, la suma de las componentes de cada columna debe ser exactamente igual a $1$: $\sum_{i=1}^n P_{ij} = 1$.
% \end{itemize}

% Construyendo la matriz $P \in \mathbb{R}^{4 \times 4}$ para nuestro grafo:
% \[
%     P = \begin{pmatrix}
%         0   & 0 & 1/3 & 0   \\
%         1/2 & 0 & 0   & 1/2 \\
%         1/2 & 0 & 1/3 & 1/2 \\
%         0   & 1 & 1/3 & 0
%     \end{pmatrix}
% \]

% \subsection*{4. Evolución Temporal y Estado Límite}
% Iterando recursivamente el proceso a partir de una distribución inicial
% $v^{(0)}$:
% \[
%     v^{(1)} = P v^{(0)}, \quad v^{(2)} = P v^{(1)} = P^2 v^{(0)}, \quad \dots \quad v^{(k)} = P^k v^{(0)}
% \]
% Si el proceso converge a una distribución estable a largo plazo ($v^{(\infty)}
%     = \lim_{k \to \infty} v^{(k)}$), este vector límite debe permanecer invariable
% al aplicar la transición:
% \[
%     v^{(\infty)} = P v^{(\infty)} \iff P v^{(\infty)} = 1 \cdot v^{(\infty)}
% \]
% Esto demuestra que el vector de ranking final de las páginas web no es otra
% cosa que un \textbf{autovector de $P$ asociado al autovalor $\lambda =
%         1$}. Toda la teoría de matrices estocásticas, el método de la potencia
% y la diagonalización que desarrollamos a continuación responden de manera
% rigurosa a:
% \begin{itemize}
%     \item ¿Por qué siempre existe el autovalor $\lambda = 1$?
%     \item ¿Por qué ningún autovalor puede tener módulo mayor a 1 ($|\lambda| \le 1$)?
%     \item ¿Bajo qué condiciones sobre el grafo el vector límite $v^{(\infty)}$ existe y es único independiente del estado inicial?
% \end{itemize}

% =========================================================================
\section{Definición del Proceso y Matrices de Markov}
% =========================================================================

\begin{definicion}[Matriz de Markov / Estocástica por Columnas]
    Una matriz cuadrada $M = (m_{ij}) \in \mathbb{R}^{n \times n}$ se dice \textbf{matriz de Markov} (o estocástica) si cumple simultáneamente:
    \begin{enumerate}
        \item \textbf{No negatividad:} $m_{ij} \ge 0$ para todo $1 \le i, j \le n$.
        \item \textbf{Suma unitaria por columnas:} Para cada columna $j \in \{1, \dots, n\}$,
            \[
                \sum_{i=1}^n m_{ij} = m_{1j} + m_{2j} + \cdots + m_{nj} = 1.
            \]
    \end{enumerate}
    En notación matricial compacta con $\mathbf{1} = {(1, 1, \ldots, 1)}^{T}$, la condición de columnas equivale a $\mathbf{1}^T M = \mathbf{1}^T$.
\end{definicion}

\begin{observacion}
    Si A es \textbf{matriz de Markov} entonces todas las entradas $a_{ij}$ $\in [0,1]$.
\end{observacion}

\begin{definicion} 
    Un proceso dado iterativamente por $v^{(k+1)} = A v^{(k)}$ con $k \geq 0$ se llama \textbf{proceso de Markov}. A cada vector $v^{(k)}$ se lo llama $\textbf{estado}$ del sistema en tiempo $k$. A la matriz $A$ se la llama $\textbf{matriz de transición}$.
\end{definicion}

% =========================================================================
\section{Autovalor de las Matrices de Markov}
% =========================================================================

\begin{proposicion}
    Sea $A \in \mathbb{R}^{n \times n}$ $\textbf{matriz de Markov}$, entonces
    \begin{enumerate}
        \item \textbf{Existencia del autovalor 1:} $\lambda = 1$ es siempre autovalor de $A$.
        \item Si $\lambda$ es $\textbf{autovalor}$ de $A$ $\rightarrow$ $\vert\vert \lambda \vert\vert$ $\leq$ 1. 
    \end{enumerate}
\end{proposicion}

\begin{demostracion}
    \begin{itemize}
        \item 1.  
            \[
                A 
                = 
                \begin{pmatrix}
                    C_1(A) | \cdots | C_n (A) 
                \end{pmatrix}
            \] 
            \[
                \rightarrow A^t = 
                    \left(
                        \begin{array}{c}
                            C_1(A) \\
                            \midrule
                            \vdots \\
                            \midrule
                            C_n(A)
                        \end{array}
                    \right) \text{los elementos de cada fila suman $1$}
            \] 
            \[
                A^t 
                \left(
                    \begin{array}{c}
                        1 \\
                        \midrule
                        \vdots \\
                        \midrule
                        1
                    \end{array}
                \right)
                =
                \left(
                    \begin{array}{c}
                        1 \\
                        \midrule
                        \vdots \\
                        \midrule
                        1
                    \end{array}
                \right)
                \rightarrow \text{$1$ es autovalor de $A^t$ con }
                \left(
                    \begin{array}{c}
                        1 \\
                        \midrule
                        \vdots \\
                        \midrule
                        1
                    \end{array}
                \right)
                \text{autovector asociado}
            \]
            \[
                A^t.v = 1 . v
            \]
            Dado que $A$ y $A^T$ comparten el mismo polinomio característico ($\det(A^T - \lambda I) = \det({(A - \lambda I)}^T) = \det(A - \lambda I)$), tienen exactamente los mismos autovalores, concluyendo que $\lambda = 1$ es autovalor de $A$.
        \item 2. \textbf{Cota espectral / Radio espectral $\lambda \le 1$:} Para todo autovalor $\lambda$ de $A$, se verifica:
            \[
                |\lambda| \le 1.
            \]
            Sea $v \neq \mathbf{0}$ autovector tal que $Mv = \lambda v$. 
            \[
                v = \left( \begin{array}{c}
                        x_1 \\
                        \midrule
                        \vdots \\
                        \midrule
                        x_n
                \end{array}
                \right)
                \hspace*{2cm}
                A \left( \begin{array}{c}
                        x_1 \\
                        \midrule
                        \vdots \\
                        \midrule
                        x_n
                \end{array}
                \right)
                = 
                \left( \begin{array}{c}
                        x_1 \\
                        \midrule
                        \vdots \\
                        \midrule
                        x_i \\
                        \midrule
                        \vdots \\
                        \midrule
                        x_n
                \end{array}
                \right)
            \]
            \[
                \lambda x_i = F_i(A) X = \sum_{j=1}^{n} a_{ij} x_j
            \]
            Tomando módulo y desigualdad triangular para cada componente $i$:
            \[
                |\lambda x_i| \leq |\lambda| |x_i| = |\sum_{j=1}^{n} a_{ij} x_j| = \sum_{j=1}^{n} a_{ij} |x_j| \hspace{1cm} \text{(pues $a_{ij}$ $\geq$ $0$, A es matriz de Markov)}
            \]
            \[
                \rightarrow 
                    \sum_{i=1}^{n} |\lambda||x_i| 
                \leq 
                    \sum_{i=1}^{n} \sum_{j=1}^{n} a_{ij} |x_j| 
                = 
                    \sum_{j=1}^{n} \sum_{i=1}^{n} a_{ij} |x_j| 
                = 
                    \sum_{j=1}^{n} \left( \sum_{i=1}^{n} a_{ij} \right) |x_j|  
                = 
                    \sum_{j=1}^{n} \left( 1 \right) |x_j|  
            \] 
                \[
                \rightarrow 
                    \sum_{i=1}^{n} |\lambda||x_i| 
                =
                    |\lambda| \sum_{i=1}^{n} |x_i| 
                \leq 
                    \sum_{j=1}^{n} |x_j|  
            \]
            Como $v \neq \mathbf{0}$, $\|v\|_1 > 0$, dividiendo resulta $|\lambda| \le 1$.
            \[
                \rightarrow
                    |\lambda| \frac{\sum_{i=1}^{n} |x_i|}{\sum_{i=1}^{n} |x_i|}
                = 
                    \frac{\sum_{j=1}^{n} |x_i|}{\sum_{i=1}^{n} |x_i|}
            \]
            \[
                \rightarrow
                    |\lambda|
                \leq
                    \frac{\sum_{j=1}^{n} |x_i|}{\sum_{i=1}^{n} |x_i|} \text{  (j, i recorren todo el vector X)}
            \]
                \[
                \rightarrow
                    |\lambda|
                \leq
                    1 
            \]
    \end{itemize}
\end{demostracion}

\begin{proposicion}
    Sea $A$ una $\textbf{matriz de Markov}$ y $\lambda$ es autovalor de $A$, $\lambda \neq$ 1. 
    \newline
    Si $v = \left( 
                \begin{array}{c}
                    x_1 \\
                    \midrule
                    \vdots \\
                    \midrule
                    x_n
                \end{array}
            \right)$ es un autovector asociado a $\lambda \rightarrow x_1+\cdots+x_n = 0 $    
\end{proposicion}

\begin{demostracion}
    Multiplicando la ecuación $Mv = \lambda v$ a izquierda por $\mathbf{1}^T$:
    \[
        \mathbf{1}^T (Mv) = (\mathbf{1}^T M)v = \mathbf{1}^T v = \lambda (\mathbf{1}^T v) \implies (1 - \lambda)(\mathbf{1}^T v) = 0.
    \]
    Como $\lambda \neq 1$, se deduce inmediatamente que $\mathbf{1}^T v = 0$.    
\end{demostracion}

% =========================================================================
\section{Estados de Equilibrio y Convergencia a Largo Plazo}
% =========================================================================

\begin{definicion}[Estado de Equilibrio / Estacionario]
    A \textbf{matriz de Markov}, $v$ es un \textbf{estado de equilibrio} si:
    \[
        A v = v \text{  ($v$ es autovector asociado al autovalor 1)}
    \]
    \begin{itemize}
        \item Con coordenadas NO negativas
        \item Con coordenadas que suman 1
    \end{itemize}
\end{definicion}
\begin{observacion}
    Si $v = \left( 
                \begin{array}{c}
                    x_1 \\
                    \midrule
                    \vdots \\
                    \midrule
                    x_n
                \end{array}
            \right)$ es un autovector asociado a $1$ 
    \[
        \rightarrow \frac{v}{\sum_{i=1}^{n}x_i} \text{  es autovector del autovalor 1 y sus coordenadas suman 1}
    \]
\end{observacion}
\begin{propiedades} $A$ \textbf{matriz de Markov}
    \begin{enumerate}
        \item Existe al menos 1 equilibrio
        \item \textbf{Propiedad clave del autoespacio $E_1$:}
            El subespacio propio $E_1 = \operatorname{Nu}(M - I)$ tiene dimensión al menos 1.
        \item No toda \textbf{matriz de Markov} es \textbf{diagonalizable}
    \end{enumerate}
\end{propiedades}



\subsection*{Descomposición Espectral y Límite $M^k$}
Supongamos que $M$ es diagonalizable: $M = C D C^{-1}$, donde $C = (w_1 \mid
    w_2 \mid \dots \mid w_n)$ contiene los autovectores y $D =
    \operatorname{diag}(\lambda_1, \dots, \lambda_n)$ con $1 = \lambda_1 \ge
    |\lambda_2| \ge \dots \ge |\lambda_n|$.

Para cualquier estado inicial $\pi^{(0)} \in \mathcal{P}_n$, expresado en la
base de autovectores $\pi^{(0)} = a_1 w_1 + \cdots + a_n w_n$:
\[
    \pi^{(k)} = M^k \pi^{(0)} = a_1 \lambda_1^k w_1 + a_2 \lambda_2^k w_2 + \cdots + a_n \lambda_n^k w_n.
\]
\begin{itemize}
    \item Si $|\lambda_i| < 1$, entonces $\lim_{k \to \infty} \lambda_i^k = 0$.
    \item Si $\lambda = 1$ es el \textbf{único} autovalor de módulo 1 con multiplicidad 1
          (autovalor dominante estricto):
          \[
              \lim_{k \to \infty} \pi^{(k)} = a_1 w_1 = \pi^* \quad \text{independientemente del estado inicial } \pi^{(0)}.
          \]
    \item En tal caso, la matriz límite $M^\infty = \lim_{k \to \infty} M^k$ existe y
          todas sus columnas son idénticas e iguales al vector de equilibrio
          estacionario:
          \[
              M^\infty = \begin{pmatrix} \pi^* & \pi^* & \cdots & \pi^* \end{pmatrix}.
          \]
\end{itemize}

\textbf{¿Cuándo NO converge el proceso?}
Si la matriz tiene otros autovalores de módulo 1 (por ejemplo, $\lambda = -1$ o raíces de la unidad $e^{i\theta}$ con matrices periódicas o bloques cíclicos), el término $\lambda_i^k$ oscilará y el límite $\lim_{k \to \infty} \pi^{(k)}$ no existirá para ciertos estados iniciales.

% =========================================================================
\section{Conexiones con Otros Temas de ALC}
% =========================================================================

\begin{table}[h!]
    \centering
    \small
    \begin{tabular}{lp{12cm}}
        \toprule
        \textbf{Tema de la Materia}          & \textbf{Conexión con Procesos de Markov}                                                                                                                                                                                                                                                                             \\
        \midrule
        \textbf{Normas Matriciales}          & $\|M\|_1 = \max_j \sum_i |m_{ij}| = 1$. Como el radio espectral verifica $\rho(M) \le \|M\|$ para toda norma inducida, esto prueba analíticamente en una sola línea que $\rho(M) \le 1$.                                                                                                                             \\
        \midrule
        \textbf{Círculos de Gershgorin}      & Para $M^T$, cada centro es $m_{jj}$ y el radio es $R_j = \sum_{i \neq j} m_{ij} = 1 - m_{jj}$. Cada disco de Gershgorin es $D(m_{jj}, 1 - m_{jj}) \subset \{z \in \mathbb{C} : |z| \le 1\}$ y toca tangencialmente al punto $z = 1$.                                                                                 \\
        \midrule
        \textbf{Método de la Potencia}       & La iteración de Markov $\pi^{(k)} = M \pi^{(k-1)}$ es \emph{exactamente} el método de la potencia aplicado a $M$ para converger al autovector dominante asociado a $\lambda_1 = 1$, con la ventaja de que no se requiere renormalizar por norma euclídea en cada paso porque la norma-1 se preserva automáticamente. \\
        \midrule
        \textbf{Teorema de Perron-Frobenius} & Para matrices regulares (o ergódicas, donde alguna potencia $M^k > 0$ es estrictamente positiva), el autovalor 1 es simple, todo otro autovalor cumple $|\lambda| < 1$, y el autovector estacionario $\pi^*$ tiene todas sus componentes estrictamente positivas ($\pi^*_i > 0$).                                    \\
        \midrule
        \textbf{Google PageRank}             & El algoritmo PageRank construye una matriz de Markov $M = d P + \frac{1-d}{n} \mathbf{1} \mathbf{1}^T$ con factor de amortiguación $d \in (0, 1)$, garantizando que la matriz sea estrictamente positiva y tenga un único vector de equilibrio calculable mediante iteraciones.                                      \\
        \bottomrule
    \end{tabular}
\end{table}

% =========================================================================
\section{Guía Práctica: Paso a Paso para Resolver Ejercicios}
% =========================================================================

\begin{enumerate}[label=\textbf{Paso \arabic*:}]
    \item \textbf{Construir la Matriz de Transición $M$:}
          \begin{itemize}
              \item Identificar los estados $1, 2, \dots, n$.
              \item Definir cada columna $j$: ubicar en la posición $(i, j)$ la probabilidad o tasa
                    de pasaje del estado $j$ al estado $i$.
              \item \emph{Checklist obligatorio:} verificar que $m_{ij} \ge 0$ y que la suma de cada columna sea exactamente 1 ($\sum_i m_{ij} = 1$).
          \end{itemize}

    \item \textbf{Hallar los Estados de Equilibrio:}
          \begin{itemize}
              \item Plantear el sistema homogéneo $(M - I)\pi = \mathbf{0}$.
              \item Triangular la matriz $(M - I)$ mediante eliminación gaussiana para encontrar el
                    núcleo $\operatorname{Nu}(M - I)$. (Recuerda que $\det(M - I) = 0$ siempre, por
                    lo que al menos una fila se anulará).
              \item Parametrizar las soluciones y normalizar el vector dividiendo por la suma de
                    sus componentes: $\pi^* = \frac{v}{\sum_i v_i}$ para que $\pi^* \in
                        \mathcal{P}_n$.
          \end{itemize}

    \item \textbf{Analizar la Existencia del Límite $\lim_{k \to \infty} \pi^{(k)}$:}
          \begin{itemize}
              \item Calcular las raíces del polinomio característico $\det(M - \lambda I) = 0$.
              \item Si $\lambda = 1$ tiene multiplicidad 1 y todos los demás autovalores cumplen
                    $|\lambda_i| < 1$, el proceso converge siempre al único $\pi^*$ para cualquier
                    $\pi^{(0)}$.
              \item Si existen autovalores complejos con $|\lambda| = 1$ o $\lambda = -1$, el
                    límite dependerá del estado inicial (puede ciclar u oscilar).
          \end{itemize}

    \item \textbf{Cálculo en $k$ Pasos Finitos o Potencias Altas:}
          \begin{itemize}
              \item Si $k$ es pequeño (ej. $k=2$ o $4$): multiplicar recursivamente $\pi^{(1)} = M
                        \pi^{(0)}$, $\pi^{(2)} = M \pi^{(1)}$, etc.
              \item Si $k$ es general: diagonalizar $M = C D C^{-1}$, calcular $M^k = C D^k C^{-1}$
                    y evaluar $\pi^{(k)} = M^k \pi^{(0)}$.
          \end{itemize}
\end{enumerate}

\end{document}